\section{Introduction}
\label{sec:intro}

Step-indexed reasoning has become the standard semantic technique
for verifying higher-order programming languages with mutable state,
recursive types, and concurrency.  From Appel and McAllester's
original step-indexed model of recursive types
\cite{appel-mcallester2001} to the modern Iris
ecosystem~\cite{jung-et-al-2018}, the same basic idea recurs: instead
of asking whether a program satisfies a specification, ask whether it
satisfies the specification at every finite observation depth $n$,
then recover the unindexed statement as the limit.  Each level of
the hierarchy is defined without circularity by appeal only to lower
levels, and a ``later'' modality $\triangleright$ mediates the
descent from one level to the next.

What the step-indexed literature does not usually say is that this
technique is, in classical set-theoretic terms, \emph{forcing}.
Cohen introduced forcing in 1963 to prove the independence of the
continuum hypothesis~\cite{cohen1963-independence-i, cohen1966}.  In
its original presentation, called \emph{ramified forcing} after
the ramified type theory of Russell and
Whitehead~\cite{russell-whitehead-PM}, the construction of the
generic extension is stratified by a well-founded hierarchy of
levels, with recursion at level $\alpha$ permitted to consult only
levels strictly below $\alpha$.  Shoenfield's 1971
reformulation~\cite{shoenfield1971} eliminated the ramification for
the set-theoretic application: the modern textbook presentation of
forcing does not use it.  But that structure is exactly the structure
of step-indexed logic: well-founded stratification, level descent, and
a forcing relation $p \Vdash \varphi$ monotone in the strength of $p$.  The forcing conditions are the natural numbers
$\omega$ with $\le$, the relation $n \Vdash \varphi$ is monotone in
$n$, the modality $\triangleright$ is the level descent operator,
and Löb's rule is well-founded induction on the forcing order,
internalized into the logic.

\paragraph{Contributions.}
We give two technical observations about step-indexed Hoare logic,
together with the forcing-style framework that makes them stateable.
All claims are mechanized in Rocq.

\begin{enumerate}

\item \textbf{A structural diagnostic for the internal Löb rule}
(Section~\ref{sec:imp-vs-lam}).  Recursive Hoare-style rules split by
the \emph{recursion locus} of the soundness predicate.  For IMP's
\texttt{while} it recurses only in the step index, where outer
induction on $n$ and internal Löb give structurally identical proofs
(\texttt{hoare\_while}, \texttt{hoare\_while\_iLob}).  For \texttt{fix}
it recurses in the program itself, where only internal Löb plays to
that structure.

\item \textbf{A characterization of L\"ob's rule by well-foundedness}
(Section~\ref{sec:lob-essential}).  In any forcing-style pre-structure
satisfying our order-theoretic axioms, L\"ob's rule holds for every
iProp iff the strict refinement relation is well-founded.  The
backward direction uses the \emph{forcing-accessibility} predicate
$\varphi_{\mathrm{acc}}\, p := \mathtt{Acc}\, \mathtt{lt}\, p$ as an
internal iProp witness forced everywhere, and a $\mathbb{Z}$ instance
where L\"ob fails on $\bot$ realises the other direction. This iff, a
step-indexed counterpart of a classical Gödel--Löb
fact~\cite{smorynski1985,boolos1993}, does not appear in the
verification literature we surveyed.

\item \textbf{A mechanized forcing-style re-presentation of
step-indexed Hoare logic} (Sections~\ref{sec:prop}--\ref{sec:lam})
with an \textbf{abstract framework parametric in the forcing notion}
(Section~\ref{sec:abstract}), instantiated at $(\omega, \le)$, the
pointwise product, and the lexicographic $\omega^2$ of transfinite
step-indexing~\cite{svendsen2016transfinite,spies2021transfinite},
which a mechanized propositional sentence \emph{separates} from
$\omega$.  Löb's rule is proved uniformly by
\texttt{well\_founded\_ind} on the strict refinement relation.

\item \textbf{The propositional fragment of the
Jaber--Tabareau--Sozeau forcing translation of
CIC~\cite{jaber-tabareau-sozeau-2012}, mechanized}
(Appendix~\ref{sec:bridge}).  We give it as a standalone
\texttt{Fixpoint} on a deep-embedded language with $\triangleright$
and prove that over an \emph{arbitrary} forcing structure it
coincides pointwise with the framework's semantics
(\texttt{jts\_is\_interp}) and validates L\"ob's rule
(\texttt{jts\_Lob}).  The \emph{term-level} translation of full CIC
remains future work.

\end{enumerate}

\paragraph{What is and is not novel.}
The re-presentation itself is not novel.  The topos-of-trees account
of step-indexing~\cite{birkedal-et-al-2012} is implicitly a
forcing-style semantics (Section~\ref{sec:topos}). The Iris
ecosystem uses internal Löb without referring to it as forcing. The
JTS translation produces, at the propositional level, structures
that coincide with our \texttt{iProp}.  The framework we give is the
smallest such account that supports the two technical observations
above.  Its purpose is to expose the structure inside these existing
constructions that makes the diagnostic and the counterexample
stateable in the first place.

\paragraph{Scope.}
We work at the level of plain Hoare triples and a partial-correctness
wp predicate, with no separation logic, heap, or concurrency.  The
$\lambda$-calculus of Section~\ref{sec:lam} is untyped, and our
recursive spec takes a single predicate $S$ as both pre- and
post-condition (the rule generalizes to distinct pre/post in the
formalization).

\paragraph{Outline.}
Section~\ref{sec:prop} gives the propositional layer.
Section~\ref{sec:imp} develops Hoare logic for IMP and proves the
\texttt{while} rule by meta-level induction.
Section~\ref{sec:lam} develops Hoare logic for an untyped
$\lambda$-calculus with \texttt{fix} and proves the soundness rule
via the internal Löb rule.  Section~\ref{sec:imp-vs-lam} articulates
the diagnostic.  Section~\ref{sec:abstract} gives the abstract
framework.  Section~\ref{sec:lob-essential} characterizes L\"ob's rule
in terms of well-foundedness of the forcing notion.
Section~\ref{sec:related} positions the contribution against the
literature, and Appendix~\ref{sec:bridge} formalizes the JTS
embedding.  The complete development is mechanized in
Rocq with no external dependencies.
